\begin{afterword}
\summaryhead

The post argues that the best evidence on a question is the evidence closest to it. When the Wright brothers claim their plane will fly, comparing their credentials with Lord Kelvin's favors Kelvin. Checking both sides' calculations makes credentials matter less, and watching the plane fly makes the calculations matter less still.

The author puts this in terms of causal graphs: evidence close to the question ``screens off'' more distant evidence, and a theorem guarantees that distant nodes cannot carry more information than closer nodes that screen them off. The same applies to judging a speaker's rationality, which is also a step removed from the question. Accusing someone of a bias does not settle a factual issue.

The author grants that we often cannot check the calculations ourselves, for lack of background, access or time. Then we must judge the speaker, but should do so with a sense that something is missing. The post ends by urging the reader to get as close to the question as possible.

\responsehead

The post states a familiar rule of evidence, and shows it working only in the case where no one needs it. When you can watch the plane fly, you do not need a principle to ignore Lord Kelvin. The rule matters when you cannot check, and there the post gives up. It concedes a ``sadly large number of times'' when the reader must judge the speaker, and offers for those times a ``hollow feeling in your heart.'' A reader facing a question in medicine or climate science, which is the normal case for a non-expert, leaves with a slogan and no method.

The example is staged, and the history it borrows undercuts it. The post cites no statement by Kelvin about the Wrights. Its ranking places calculations above authority, but the Wrights' calculations are what failed. Their 1901 glider gave about a third of the lift they had calculated, because the calculations used a coefficient that had been accepted for more than a century and was wrong. Checking their arithmetic would not have helped. They found the error by building a wind tunnel. Calculations inherit the authority of their inputs, and the post's picture of a clean progression from credentials to proof hides that.

The theory is overstated. The theorem the post invokes is the data-processing inequality, and it says something narrower than the post's first claim. The guarantee holds only when the closer evidence screens the distant evidence off, and it says nothing about reading the closer evidence wrongly. What a reader observes is never the calculation itself but their own reading of it. The post asks the reader to accept ``a theorem of these causal graphs'' without a name or a source, which is an argument from authority in a post against arguments from authority. It then quotes a fiction-writing manual for support.

The footnote on the exception is the most revealing sentence in the post. Where the author concedes that claims sometimes cannot be checked, he sends the reader to his own defence of molecular nanotechnology, a claim he supports and which, by his own account there, is ``merely a best guess'' that no experiment had tested. I infer that the exception is there for the author's causes, and that the rule is for everyone else's.

In short: a rule of evidence shown only where it is obvious, abandoned where it is needed, and illustrated with a history that undercuts it.
\end{afterword}
